\section{\label{sec:3} Results}

\begin{figure}[t]
\includegraphics[width=17cm]{../figures/Figure4}
\caption{\label{fig:FIG4}{Time course of vertical localization performance. White circles indicate localization with free ears, black circles and grey squares indicate localization with the first and second earmolds respectively. Data points represent the mean across participants (mold 1: n = 15, mold 2: n = 12). Error bars show the standard error of the mean. Individual learning trajectories are shown in light grey. \textbf{A} Time course of elevation gain, \textbf{B} vertical localization error and \textbf{C} vertical response variability (SD) throughout the two consecutive adaptation periods. Participants elevation gain was reduced while localization error and variability were increased after insertion of earmolds. Elevation gain and vertical error recovered within each of the two five-day adaptation periods, while response variability further increased. Participants' gain recovered faster with the first molds than with the second set. After each mold removal, elevation gain returned to baseline levels whereas vertical error increased compared to the first test with free ears (white dots on day 5 and 10). Although elevation gain with the first molds was slightly decreased after participants adaptated to the second set, the persistence of adaptation did not differ between molds.}}
\end{figure}

\subsection{Earmolds reduced vertical localization performance}